\begin{frame}{핵심 아이디어: 함수근사와 일반화}
  \begin{itemize}
    \item $Q^{*}$는 상태-행동마다 다른 값을 주는 함수 ---
      표 대신 \textbf{수식으로} 표현할 함수족 $Q(s,a;\theta)$ ($\theta$ = 가중치)로
      근사: \kb{함수근사}(function approximation).
    \item 표: ``같은 칸'' 상태끼리만 정보 공유. 다른 칸은 아무리 비슷해도
      배운 것을 \textbf{아무것도} 전달하지 못한다.
  \end{itemize}
  \vskip0.5em
  \begin{block}{\kb{일반화}(generalization)}
    함수가 매끄러우면, \textbf{한 번 본 상태}에서 배운 것이
    \textbf{한 번도 본 적 없는 비슷한 상태}까지 자동으로 퍼져나간다.
    ML1의 ``학습 데이터에 없었던 새 입력에 대한 예측''과 같은 개념.
  \end{block}
  \vskip0.4em
  \kq{함수근사의 진짜 이유는 메모리가 아니라 \textbf{일반화}다.}
\end{frame}
